\section{Experiments}\label{sec:experiments}

\subsection{Generative Features}\label{sec:generative_features}

\cite{MelodyTranscriptionViaGenerativePreTraining} use generative features extracted from Jukebox~\cite{Jukebox} to improve performance for melody transcription. They also produce a chord transcription model using the same methodology but do not report results. We test generative features using MusicGen~\cite{MusicGen} as a feature extractor. This was for several reasons. MusicGen is a newer model. It has several different sizes of models that can be tested to ablate on the complexity of the model. It has a fine-tuned variant called MusiConGen~\cite{MusiConGen} which is used for synthetic data generation in Section~\ref{sec:synthetic_data}. All model weights are available on the HuggingFace Hub.\footnote{\url{https://huggingface.co/docs/hub/en/index}} Finally, its training data was properly licensed, unlike Jukebox. We leave the results using Jukebox features and other models for future work. Details of the feature extraction process can be found in Appendix~\ref{app:generative_feature_extraction}.

% Overlapping $5$ second chunks of audio were fed through MusicGen in a batched fashion. This first requires passing the audio through the pre-trained Encodec audio tokeniser~\cite{Encodec}. These are then fed through the language model. We take the output logits as the representation for each frame. The model outputs logits in four `codebooks', each $2048$-dimensional vectors, intended to represent different granularities of detail in the audio. Audio segments are overlapped such that every frame has context from both directions. The multiple representations for each frame are averaged. Finally, these representations are upsampled. The model operates at a frame rate of 50Hz. To compute a representation with the same frame length as the CQT, we take the mean over the frames outputted by the model closest to the centre of the CQT frame. In case averaging over frames dampened the signal, we also tried linearly interpolating between the two closest frames outputted by the model. However, this was empirically found to perform slightly worse. Results can be found in Appendix~\ref{app:linear_interpolation_vs_area_averaging}. This feature extraction required the use of NVIDIA RTX A6000 GPUs. The extraction process took approximately 4 hours for each model over the entire dataset. Due to computational constraints, pitch augmentation was removed for generative feature experiments.

As the representations are $2048$-dimensional, it is computationally infeasible to feed this directly into the GRU. Instead, using fully connected layers, we project these vectors down to a power of $2$, from $16$ to $1024$. The best representation is with $64$ dimensions, although results show no clear trend. Results are left to Appendix~\ref{app:projection_dimensionality}.

We test different variants of MusicGen, including musicgen-large (3.3B parameters), musicgen-small (330M), musicgen-melody (1.5B) and MusiConGen (1.5B). We also test different reductions of the four `codebooks' outputted by the language model. This includes taking each codebook on its own, the element-wise mean across codebooks and the concatenation of all four. For further details on codebooks, see the work of \cite{MusicGen}. The best model was found to be musicgen-large and the best reduction to average over the four codebooks. However, results were close and are left to Appendices \ref{app:generative_feature_extraction_models} and \ref{app:generative_feature_extraction_reductions} for lack of interest.

% It might be argued that it is better to use musicgen-small to save on computational cost but the dimensionality of the codebooks is the same across all model variants. Once the features have been extracted, they are all equally as expensive to train on.

% Surprisingly, the concatenated representation performs worse than the averaged representation as it contains at least as much information. However, if the information provided by each codebook is essentially the same, then there are no reasons that the concatenated representation should perform better and training with Adam may simply find a worse minimum. Training on $8192$-dimensional vectors is also computationally expensive.

To test whether or not these features help when compared with a CQT, we test with the CQT only, generative features only and a concatenation of the two. The results are shown in Table~\ref{tab:gen_feature_comparison}. Although the generative features perform worse than the CQT, they contain information useful for chord recognition with an accuracy of $58.7\%$. Performance remains largely the same when the CQT and generative features are used together. This experiment was run multiple times, with similar results each time. There is no clear evidence that the generative features from MusicGen provide any benefit over just using the CQT. 

This conclusion is surprising as \cite{MelodyTranscriptionViaGenerativePreTraining} claim that generative features are better than hand-crafted features for the related task of melody recognition. However, they only compare to mel-spectrograms, which may not perform as well as CQTs, as they certainly do not for chord recognition. Observations here cast doubt on how well their claims generalise to chord recognition. Features extracted from other generative models such as Jukebox~\cite{Jukebox} or MusicLM~\cite{MusicLM} may perform better as MusicGen was trained with only licensed music, severely limiting the portion of the full distribution it learnt structure about. The comparison is left for future work. Another explanation is that \cite{MelodyTranscriptionViaGenerativePreTraining} use Mel-spectrogram features as baselines. We perform a further experiment with results in Appendix~\ref{app:spectrogram_variants} showing that CQTs outperform Mel-spectrograms for ACR. Thus the CQT baseline may be stronger in ACR than Mel-spectrograms for melody transcription.

Given the lack of improvement and the drastically increased computational cost associated with extracting features and training the model, we do not proceed with training on generative features.

\begin{table}
    \centering
    \begin{tabular}{lccccc}
        \toprule
        features & accuracy & root  & third & seventh & mirex \\  
        \midrule
        gen$\cdot$CQT  & \textbf{61.0}     & \textbf{80.3}  & \textbf{77.0}  & \textbf{63.3}    & 78.4  \\
        gen      & 58.7     & 77.6  & 74.3  & 60.9    & 77.5  \\
        CQT      & \textbf{61.0}     & 79.8  & 76.8  & 63.2    & \textbf{79.3}  \\
        \bottomrule
    \end{tabular}
    \caption{Comparison of CQT and generative features with a concatenation of the two, denoted as gen$\cdot$CQT. The concatenation performs the best on most metrics but not by enough to claim is meaningfully better. }\label{tab:gen_feature_comparison}
\end{table}


\subsection{Synthetic Data Generation}\label{sec:synthetic_data}

Given the success of pitch augmentation for ACR, it is sensible to look for other sources of data. Further augmentation is possible by adding noise and time-stretching. However, these do not provide new harmonic structures or create new instances of rare chords. Instead, we look to generate new data, taking into account our understanding of harmonic structure. Generation would be possible through automatic arrangement, production and synthesis software. However, this is a complex task requiring a lot of human input and is unlikely to produce a sufficient variety of timbres, instrumentation and arrangement. Instead, we use a recent chord-conditioned generative model called MusiConGen~\cite{MusiConGen}. We use this over CocoMulla~\cite{CocoMulla} as the method of chord-conditioning has a more straightforward interface, doesn't require reference audio and the authors claim that it adheres more closely to its chord conditions. Indeed, they feed the outputted audio through BTC~\cite{BTC} and find a \texttt{triads} score of $71\%$ using the chord conditions as the ground truth. While far from perfect, it suggests the model can generate audio that mostly adheres to the chord conditions.

We generate $1210$ songs, each $30$ seconds long, to mimic the size of the \emph{pop} dataset. We refer to this dataset as \emph{synth}. This is split into a train/validation/test split in the same fashion as for the \emph{pop} dataset. Generating a larger order of magnitude of audio would require more computational resources. While the model supports auto-regressive generation for longer audio, its outputs become incoherent using the provided generation functions. It also sometimes produces incoherent output with $30$ seconds of generation, but this was much less common. The details of generating the synthetic dataset can be found in Appendix~\ref{app:synthetic_data_gen}. 

The process generates a very different chord distribution than the one in \emph{pop} with many more instances of upper extensions and rare qualities. This is intended to provide the model with many more instances of rare chords. To offset this distribution shift, we calibrate the probabilities outputted by the model. To encourage root-invariance, calibration, terms are averaged over roots for each quality. The details of calculating calibration terms are described in Appendix~\ref{app:calibration}.

We manually inspect outputs from MusiConGen. In general, the outputs are good. They consistently stick to the specified BPM and usually stick to the chord conditions. Outputs are occasionally musically strange with jarring drum transients and unrealistic chord transitions. They also did not have an enormous variety in timbre and instrumentation. Nonetheless, most examples have sensible annotations that one would expect a human to be able to annotate well. 

We empirically compare a model trained on only \emph{pop}, trained on only synthetic data and trained on both. While the latter results in more training data per epoch, convergence is always reached, so these are fair comparisons. We test the models on the \emph{pop} and synthetic data validation splits. Given the increased instances of rare chords in the synthetic data, we remove the weighting on the loss function for the models trained on synthetic data.

Table~\ref{tab:synthetic_data} shows the results. The model trained on both datasets performs very similarly to the model only trained on \emph{pop}. That is, the model has not overfit to the synthetic data. However, it has also not improved performance. Training on synthetic data drastically improves the accuracy on the \emph{synth} validation split. This is likely due to the unique chord distribution of the generated jazz progressions and the consistent instrumentation. The model trained on both datasets performs the best, with an accuracy of $51.0\%$.

\begin{table}[h]
    \centering
    \begin{tabular}{lcccc}
        \toprule
        train & \emph{pop} & \emph{pop} & \emph{pop} & \emph{synth} \\  
         & acc & third & seventh & acc \\
        \midrule
        \emph{pop} & \textbf{62.1} & 77.7 & 64.4 & 24.2 \\
        \emph{synth} & 48.6 & 64.4 & 50.6 & 44.8 \\
        both & \textbf{62.1} & \textbf{77.9}  & \textbf{64.6} & \textbf{51.0} \\
        \bottomrule
    \end{tabular}
    \caption{Results for models trained on \emph{pop}, \emph{synth} and both together. 
    % Unsurprisingly, models trained on \emph{pop} perform better on \emph{pop} and models trained on \emph{synth} perform better on \emph{synth}. The model trained on both does not do meaningfully better on the \emph{pop} validation split, but performance increases on the \emph{synth} validation split. Training only on \emph{pop} results in an accuracy of just $24.2\%$ on \emph{synth}, much lower than reported by \cite{MusiConGen}. However, this may be due to the unrealistic distribution of chords in the generated sequences. 
    }\label{tab:synthetic_data}
\end{table}

For further insight, the difference in confusion matrices over chord qualities is plotted in Appendix~\ref{app:cm_synthetic_data}. There are few clear trends. Recall increases for some rare qualities and decreases for others. Notably, recall on \texttt{min7} improves by $13\%$ and the model is much better at predicting the third in dominant and min7 chords. This is likely why the third and seventh accuracies increase slightly on the \emph{pop} validation set. 

A problem may be that the chord-conditioned model was trained using an ACR model~\cite{BTC} which is fundamentally limited in its representations of harmonic structure. These limitations are the same ones we are trying to break through using synthetic data. This is akin to the problem of model collapse when recursively training on synthetic data~\cite{syntheticcollapse}.

While performance does not meaningfully improve on \emph{pop}, the lack of overfitting and the gain on the synthetic data provide hope that synthetic data could prove a useful tool for ACR with further work and improved generative models. Furthermore, producing hand-annotated data is error-prone and subject to human interpretation. Synthetic data may produce non-identifiable data points but is consistent and error-free. Poor performance on synthetic data can only be explained by failures of the model or indeterminacy of the generated samples, not incorrect labels. Regardless, given the lack of performance increase, we do not continue to train on synthetic data.

\subsection{Beat Synchronisation}\label{sec:beat-synchronisation}

Chords exist in time. Musicians interpret chords in songs as lasting for a certain number of beats, not a fixed length in time. In its current form, the model outputs frame-wise predictions. While these could be stitched together to produce a predicted chord progression or beat-wise predictions could be made as a post-processing step, we decide to implement a model that outputs beat-wise predictions directly. This allows the model to use information from the entire duration of the beat to make its prediction.

Following the methodology of~\cite{MelodyTranscriptionViaGenerativePreTraining}, we first detect beats using \texttt{madmom}~\cite{madmom}. This returns a list of time steps where beats have been detected. We first verify that these beats are plausible by computing the maximum accuracy a model could attain if predicting chords at the whole beat level. This is done by iterating over each beat interval and assigning the chord with maximum overlap with the ground truth. This yields an accuracy of $97.1\%$, thus we are satisfied that the estimated beats are accurate enough to be used.

To calculate features for a beat interval, we average all CQT features whose centre is contained within the interval. The CQT is calculated using a hop length of $1024$. The shorter hop length is used to minimise the effect of CQT frames with partial overlap with two beat intervals and ensures that each beat has many CQT frames associated with it. These representations have the added benefit of decreased computational cost as beats have a lower frequency than frames.

% \{MelodyTranscriptionViaGenerativePreTraining} task the model with producing predictions for every 1/16th of a beat. However, such fast transitions are not necessary for chord recognition. The possible accuracy of $97.1\%$ over beats suggests that the model need only predict once every beat interval.

We test the model with different divisions and groupings of beats. This tests the assumption that whole beats are fine-grained enough for the model to make good predictions. We include tests where beat intervals are sub-divided into two or four, or beat intervals are joined into groups of two. We refer to this as the \emph{beat division}. We also test a \emph{oracle} beat division where the true chord transitions are taken as the beats. This is not a fair comparison as the model should not have access to true transition timings. However, it does provide an idea of how the model would fare if it could perfectly predict chord transitions. HMM smoothing is removed for beat-wise models as it is no longer necessary.

Results are shown in Table~\ref{tab:beat_division}. Using beat-wise predictions does not affect performance compared to frame-wise predictions. The model performs just as well with a beat division of $1$ as with a beat division of $1/2$ or $1/4$. The model performs worse with a beat division of $2$, though it only loses $6\%$ accuracy. This suggests that most chord transitions occur at frequencies lower than the beats produced by \texttt{madmom}.

The \emph{oracle} model performs slightly worse on accuracy metrics but attains a very high \texttt{mirex} score of $90.4\%$, the highest seen in the literature. Data from the CQT producing a \texttt{mirex} of $90.4\%$ provides hope for significant improvements in the field. Why the model performs worse on accuracy metrics is not clear. It may be because averaged features from longer time periods provide better information as to the pitch classes present but dampen signal regarding the root note. Indeed, the mean chord duration is $1.68$ seconds while a CQT frame is $0.093$ seconds. Why this effect is not observed when predicting at the beat level is also not clear. 

% Further analysis may reduce the gap between \texttt{mirex} and accuracy and improvements on ACR.

\begin{table*}[h]
    \centering
    \begin{tabular}{lccccc}
        \toprule
        beat division & acc & root & third & seventh & mirex \\  
        \midrule
        1/4                 & 62.3           & 81.3          & 78.3          & 64.5           & 79.6         \\
        1/2                 & \textbf{62.8}  & 81.5          & \textbf{78.6} & \textbf{65.1}  & 79.9         \\
        1                   & 62.3           & 81.3          & 78.1          & 64.6           & 80.0         \\
        2                   & 56.4           & 74.7          & 71.6          & 58.5           & 73.4         \\
        % 4 & 50.8           & 68.0          & 64.5          & 52.8           & 67.2         \\
        none                & 62.5           & \textbf{81.7} & 78.2          & 64.7           & 80.2         \\
        oracle             & 61.1           & 79.6          & 76.1          & 63.4           & \textbf{90.4}\\
        \bottomrule
    \end{tabular}
    \caption{Results for different beat divisions. The beat division of `none' refers to a frame-wise \emph{CRNN} with a hop length of $4096$ and HMM smoothing applied, `1' refers to the beats from \text{madmom} and other numbers are the corresponding subdivision/grouping, and `oracle' refers to the intervals that are taken from the chord labels. 
    % Note that this `perfect' model is not a fair test. Beat-wise predictions with beat intervals of $1$ and below show performance decrease compared to frame-wise predictions. Beat intervals greater than $1$ increasingly suffer from being forced to assign predictions to larger periods of time that may align with true chord transitions. Interestingly, a beat interval of 2 still performs relatively well. A notable result is the \texttt{mirex} score of the `perfect' model. A score of $90.4\%$ is the highest of any in the literature. Despite this, its accuracy does not improve. 
    }\label{tab:beat_division}
\end{table*}

\subsection{Final Results}\label{sec:test-set}

For final results, we retrain select models on the combined training and validation splits and test on the held-out test split. This is an 80/20\% train/test split. We consider the original \emph{CRNN} with no improvements, the \emph{CRNN} with a weighted and structured loss and HMM smoothing, concatenating generative features with CQTs, pitch augmentation, beat-wise predictions, `perfect' beat-wise predictions and training on synthetic data. 

Results are shown in Table~\ref{tab:test_set}. Observations are largely similar to those found previously. Weighted and structured loss with smoothing improves accuracy by $1.2\%$ and pitch shifting improves accuracy by a further $1\%$. Generative features do not help and synthetic data improves performance by an additional $0.6\%$. This alone is not a clear enough signal that training on synthetic data is better, but it provides hope for further work. Finally, beat-wise predictions maintain the same performance. The `perfect' model achieves the highest performance across all metrics. The \texttt{mirex} of $90\%$ found on the validation set has reduced to $88.7\%$, and the gap with accuracy has narrowed compared to previous results in Table~\ref{tab:beat_division}. This is still a significant result and suggests that there is hope for breaking through the `glass ceiling'. The mean class-wise accuracy does not improve past $19.7\%$ without `perfect' beats. The median in this case is $6\%$. The model's performance on the long tail remains poor. Using a greater $\alpha$ in weighting may improve the picture but would require sacrificing accuracy.

Another observation is that the model's accuracy did not improve on the test set with the additional training data. Accuracy of the \emph{CRNN} with pitch shifts trained on only $60\%$ of \emph{pop} attained an accuracy of $64.3\%$, while the model trained on the full $80\%$ training split attains $63.7\%$. The discrepancy suggests that more data from the same distribution may not improve performance. 

\begin{table*}[h]
    \centering
    \begin{tabular}{lcccccc}
        \toprule
         & acc & root & third & seventh & mirex & acc\textsubscript{class} \\
        \midrule
        \emph{CRNN}                & 61.6  & 79.3  & 76.5  & 63.3  & 80.6  & 18.9 \\
        + weighted/structured/HMM  & 62.8  & 80.9  & 78.3  & 64.6  & 80.6  & 18.6 \\
        + gen features             & 62.7  & 80.4  & 77.9  & 64.6  & \textbf{80.6}  & 19.5 \\
        + pitch shift              & 63.8  & \textbf{82.4}  & 79.4  & 65.6  & 80.3  & \textbf{19.7} \\
        + synthetic data           & \textbf{64.4}  & 82.2  & \textbf{79.9}  & \textbf{66.3}  & \textbf{81.5}  & 18.3 \\
        + beats                    & 63.7  & 82.4  & 79.7  & 65.6  & 80.9  & \textbf{19.7} \\
        \midrule
        + oracle beats             & 65.8  & 84.5  & 81.7  & 67.6  & 88.7  & 21.2 \\
        \bottomrule
    \end{tabular}
    \caption{Final results from various experimental setups on the test set. 
    % The `perfect beats' model assumes oracle beat tracking and achieves the highest results across all metrics. It is excluded when considering the best results for each metric as it is not a fair comparison. Beat-wise models do not use synthetic data. Observations echo those found previously on the validation set. Adding the weighted loss, structured loss term and an HMM smoother improves accuracy by $1.2\%$. Generative features do not help. Pitch shifting improves metrics by $1\%$ and synthetic data by a further $0.6\%$.
    }\label{tab:test_set}
\end{table*}